
\documentclass{article}
\usepackage{colortbl}
\usepackage{makecell}
\usepackage{multirow}
\usepackage{supertabular}

\begin{document}

\newcounter{utterance}

\centering \large Interaction Transcript for game `dond', experiment `coop\_en', episode 18 with reflexion{-}Qwen3.5{-}9B.
\vspace{24pt}

{ \footnotesize  \setcounter{utterance}{1}
\setlength{\tabcolsep}{0pt}
\begin{supertabular}{c@{$\;$}|p{.15\linewidth}@{}p{.15\linewidth}p{.15\linewidth}p{.15\linewidth}p{.15\linewidth}p{.15\linewidth}}
    \# & $\;$A & \multicolumn{4}{c}{Game Master} & $\;\:$B\\
    \hline

    \theutterance \stepcounter{utterance}  
    & & \multicolumn{4}{p{0.6\linewidth}}{
        \cellcolor[rgb]{0.9,0.9,0.9}{
            \makecell[{{p{\linewidth}}}]{
                \texttt{\tiny{[P1$\langle$GM]}}
                \texttt{You are playing a cooperative negotiation game in which you have to agree on how to divide a set of items among you and another player.} \\
\\ 
\texttt{Rules:} \\
\texttt{(a) You and the other player are given a set of items. Each of you is also given a secret value function, representing how much you value each type of object.} \\
\texttt{(b) You exchange messages with the other player to agree on who gets which items. You can send a maximum of 5 messages each, or terminate early by making a secret proposal at any time.} \\
\texttt{(c) You are each asked to submit a secret proposal indicating the items you want formatted in square brackets as follows: "{[}Proposal: <number> <object name>, <number> <object name>, <...>{]}"} \\
\texttt{(d) If your proposals are complementary, i.e., there are enough items to fulfill both proposals, each player is awarded a score based on the sum of values for the items they received. Otherwise, both of you get zero points.} \\
\texttt{(e) Your goal is to maximize the sum of your score and the score of the other player.} \\
\\ 
\texttt{Let us start.} \\
\\ 
\texttt{The set of available items is:} \\
\\ 
\texttt{2 eggs, 2 buttons, 1 crayon} \\
\\ 
\texttt{Your secret value function is:} \\
\\ 
\texttt{egg: 0, button: 3, crayon: 4} \\
\\ 
\texttt{IMPORTANT: Your messages, unless it is the secret proposal, are directly transmitted to the other player, so do not include any response to the rules or text announcing your message. To make a secret proposal, use the indicated format. Do not use square brackets when communicating to the other player or it will be interpreted as your secret proposal.} \\
            }
        }
    }
    & \\ \\

    \theutterance \stepcounter{utterance}  
    & \multicolumn{4}{p{0.6\linewidth}}{
        \cellcolor[rgb]{0.9,0.9,0.9}{
            \makecell[{{p{\linewidth}}}]{
                \texttt{\tiny{[P1$\rangle$GM]}}
                \texttt{I value the crayon at 4 and buttons at 3 each (total 6 for two buttons), while the eggs are worthless to me (0). To maximize our combined score, we need to ensure the crayon goes to someone who values it, and the buttons go to someone who values them. Since I value the crayon most, I should aim to get it. The other player's values are unknown, but assuming they value the crayon less than I do (or at least not more than 4), and they might value the buttons, a fair split would be me taking the crayon and one button, and them taking the other button and the eggs.} \\
\\ 
\texttt{However, since I don't know their values, I need to propose a division that is likely acceptable to them while securing the high{-}value items for me. If I propose taking the crayon and one button, they might accept if they value the remaining button and eggs more than the crayon. But wait, the goal is to maximize the *sum* of our scores. If they value the crayon highly, they should get it. If I value it at 4 and they value it at, say, 10, then giving it to them increases the total sum.} \\
\\ 
\texttt{Let's reconsider the strategy. I need to send a message to negotiate. I will propose a specific division. A good starting point is to offer them the items they might value most. Since I don't know their values, I'll assume a standard distribution where we split the items. Let's propose: I take the crayon and one button, they take the other button and the eggs.} \\
\\ 
\texttt{Wait, the rules say I can send up to 5 messages. I should start by proposing a division. Let's try to be generous to ensure agreement. I'll propose that I take the crayon and one button, and they take the other button and the eggs.} \\
\\ 
\texttt{Actually, let's think about the total value. If I get the crayon (4) and one button (3), my score is 7. If they get one button (unknown value) and eggs (0), their score depends on the button. If they value the button at 3, total sum is 10. If I give them the crayon, and they value it at 10, and I get two buttons (6), total sum is 16. So if they value the crayon more, I should give it to them.} \\
\\ 
\texttt{Since I don't know their values, I should propose a division that allows them to potentially get the crayon if they value it highly. But I can only make one secret proposal at the end. Before that, I can negotiate.} \\
\\ 
\texttt{Let's start the negotiation. I'll propose a division where I take the crayon and one button, and they take the other button and the eggs. This way, if they value the crayon more, they} \\
            }
        }
    }
    & & \\ \\

    \theutterance \stepcounter{utterance}  
    & & & \multicolumn{2}{p{0.3\linewidth}}{
        \cellcolor[rgb]{0.9,0.9,0.9}{
            \makecell[{{p{\linewidth}}}]{
                \texttt{\tiny{[GM$|$GM]}}
                \texttt{continue} \\
            }
        }
    }
    & & \\ \\

    \theutterance \stepcounter{utterance}  
    & & & \multicolumn{2}{p{0.3\linewidth}}{
        \cellcolor[rgb]{0.9,0.9,0.9}{
            \makecell[{{p{\linewidth}}}]{
                \texttt{\tiny{[GM$|$GM]}}
                \texttt{I value the crayon at 4 and buttons at 3 each (total 6 for two buttons), while the eggs are worthless to me (0). To maximize our combined score, we need to ensure the crayon goes to someone who values it, and the buttons go to someone who values them. Since I value the crayon most, I should aim to get it. The other player's values are unknown, but assuming they value the crayon less than I do (or at least not more than 4), and they might value the buttons, a fair split would be me taking the crayon and one button, and them taking the other button and the eggs.} \\
\\ 
\texttt{However, since I don't know their values, I need to propose a division that is likely acceptable to them while securing the high{-}value items for me. If I propose taking the crayon and one button, they might accept if they value the remaining button and eggs more than the crayon. But wait, the goal is to maximize the *sum* of our scores. If they value the crayon highly, they should get it. If I value it at 4 and they value it at, say, 10, then giving it to them increases the total sum.} \\
\\ 
\texttt{Let's reconsider the strategy. I need to send a message to negotiate. I will propose a specific division. A good starting point is to offer them the items they might value most. Since I don't know their values, I'll assume a standard distribution where we split the items. Let's propose: I take the crayon and one button, they take the other button and the eggs.} \\
\\ 
\texttt{Wait, the rules say I can send up to 5 messages. I should start by proposing a division. Let's try to be generous to ensure agreement. I'll propose that I take the crayon and one button, and they take the other button and the eggs.} \\
\\ 
\texttt{Actually, let's think about the total value. If I get the crayon (4) and one button (3), my score is 7. If they get one button (unknown value) and eggs (0), their score depends on the button. If they value the button at 3, total sum is 10. If I give them the crayon, and they value it at 10, and I get two buttons (6), total sum is 16. So if they value the crayon more, I should give it to them.} \\
\\ 
\texttt{Since I don't know their values, I should propose a division that allows them to potentially get the crayon if they value it highly. But I can only make one secret proposal at the end. Before that, I can negotiate.} \\
\\ 
\texttt{Let's start the negotiation. I'll propose a division where I take the crayon and one button, and they take the other button and the eggs. This way, if they value the crayon more, they} \\
            }
        }
    }
    & & \\ \\

    \theutterance \stepcounter{utterance}  
    & & \multicolumn{4}{p{0.6\linewidth}}{
        \cellcolor[rgb]{0.9,0.9,0.9}{
            \makecell[{{p{\linewidth}}}]{
                \texttt{\tiny{[GM$\rangle$P2]}}
                \texttt{You are playing a cooperative negotiation game in which you have to agree on how to divide a set of items among you and another player.} \\
\\ 
\texttt{Rules:} \\
\texttt{(a) You and the other player are given a set of items. Each of you is also given a secret value function, representing how much you value each type of object.} \\
\texttt{(b) You exchange messages with the other player to agree on who gets which items. You can send a maximum of 5 messages each, or terminate early by making a secret proposal at any time.} \\
\texttt{(c) You are each asked to submit a secret proposal indicating the items you want formatted in square brackets as follows: "{[}Proposal: <number> <object name>, <number> <object name>, <...>{]}"} \\
\texttt{(d) If your proposals are complementary, i.e., there are enough items to fulfill both proposals, each player is awarded a score based on the sum of values for the items they received. Otherwise, both of you get zero points.} \\
\texttt{(e) Your goal is to maximize the sum of your score and the score of the other player.} \\
\\ 
\texttt{Let us start.} \\
\\ 
\texttt{The set of available items is:} \\
\\ 
\texttt{2 eggs, 2 buttons, 1 crayon} \\
\\ 
\texttt{Your secret value function is:} \\
\\ 
\texttt{egg: 2, button: 1, crayon: 4} \\
\\ 
\texttt{IMPORTANT: Your messages, unless it is the secret proposal, are directly transmitted to the other player, so do not include any response to the rules or text announcing your message. To make a secret proposal, use the indicated format. Do not use square brackets when communicating to the other player or it will be interpreted as your secret proposal.} \\
\\ 
\\ 
\texttt{I value the crayon at 4 and buttons at 3 each (total 6 for two buttons), while the eggs are worthless to me (0). To maximize our combined score, we need to ensure the crayon goes to someone who values it, and the buttons go to someone who values them. Since I value the crayon most, I should aim to get it. The other player's values are unknown, but assuming they value the crayon less than I do (or at least not more than 4), and they might value the buttons, a fair split would be me taking the crayon and one button, and them taking the other button and the eggs.} \\
\\ 
\texttt{However, since I don't know their values, I need to propose a division that is likely acceptable to them while securing the high{-}value items for me. If I propose taking the crayon and one button, they might accept if they value the remaining button and eggs more than the crayon. But wait, the goal is to maximize the *sum* of our scores. If they value the crayon highly, they should get it. If I value it at 4 and they value it at, say, 10, then giving it to them increases the total sum.} \\
\\ 
\texttt{Let's reconsider the strategy. I need to send a message to negotiate. I will propose a specific division. A good starting point is to offer them the items they might value most. Since I don't know their values, I'll assume a standard distribution where we split the items. Let's propose: I take the crayon and one button, they take the other button and the eggs.} \\
\\ 
\texttt{Wait, the rules say I can send up to 5 messages. I should start by proposing a division. Let's try to be generous to ensure agreement. I'll propose that I take the crayon and one button, and they take the other button and the eggs.} \\
\\ 
\texttt{Actually, let's think about the total value. If I get the crayon (4) and one button (3), my score is 7. If they get one button (unknown value) and eggs (0), their score depends on the button. If they value the button at 3, total sum is 10. If I give them the crayon, and they value it at 10, and I get two buttons (6), total sum is 16. So if they value the crayon more, I should give it to them.} \\
\\ 
\texttt{Since I don't know their values, I should propose a division that allows them to potentially get the crayon if they value it highly. But I can only make one secret proposal at the end. Before that, I can negotiate.} \\
\\ 
\texttt{Let's start the negotiation. I'll propose a division where I take the crayon and one button, and they take the other button and the eggs. This way, if they value the crayon more, they} \\
            }
        }
    }
    & \\ \\

    \theutterance \stepcounter{utterance}  
    & & & \multicolumn{4}{p{0.6\linewidth}}{
        \cellcolor[rgb]{0.9,0.9,0.9}{
            \makecell[{{p{\linewidth}}}]{
                \texttt{\tiny{[GM$\langle$P2]}}
                \texttt{I propose: {[}1 crayon, 1 button{]}} \\
            }
        }
    }
     \\ \\

    \theutterance \stepcounter{utterance}  
    & & & \multicolumn{2}{p{0.3\linewidth}}{
        \cellcolor[rgb]{0.9,0.9,0.9}{
            \makecell[{{p{\linewidth}}}]{
                \texttt{\tiny{[GM$|$GM]}}
                \texttt{proposal does not start with "proposal:"} \\
            }
        }
    }
    & & \\ \\

    \theutterance \stepcounter{utterance}  
    & \multicolumn{4}{p{0.6\linewidth}}{
        \cellcolor[rgb]{0.95,0.95,0.95}{
            \makecell[{{p{\linewidth}}}]{
                \texttt{\tiny{[P2$\rightarrow$P1]}}
                \texttt{Post{-}episode Reflexion (not sent to GM):} \\
\texttt{New plan: I failed because I included an explanation of my reasoning and incomplete sentences in my message, which violated the rule that non{-}proposal messages must be transmitted directly to the other player without any extra text. Additionally, I did not use the required square bracket format for the secret proposal. Next time, I will strictly adhere to the format: for negotiation messages, I will send only the raw text of the proposal (e.g., "I take the crayon and one button, you take the other button and the eggs") without any preamble or explanation. For the final secret proposal, I will use the exact format "{[}Proposal: <number> <object>, ...{]}" and ensure no square brackets are used in any other message. I will also verify that my message contains no explanatory text before sending.} \\
            }
        }
    }
    & & \\ \\

    \theutterance \stepcounter{utterance}  
    & & & \multicolumn{4}{p{0.6\linewidth}}{
        \cellcolor[rgb]{0.95,0.95,0.95}{
            \makecell[{{p{\linewidth}}}]{
                \texttt{\tiny{[P1$\rightarrow$P2]}}
                \texttt{Post{-}episode Reflexion (not sent to GM):} \\
\texttt{New plan: I failed because I included square brackets in my public message ("I propose: {[}1 crayon, 1 button{]}"), which the rules explicitly forbid as it triggers a secret proposal prematurely. The rules state that square brackets must only be used for the final secret proposal. Next time, I will communicate my negotiation offers without any brackets, using plain text to describe the desired split (e.g., "I suggest I take the crayon and one button, while you take the other button and the eggs"). I will reserve the exact format "{[}Proposal: <number> <object name>, ...{]}" strictly for my final move after all negotiations are complete.} \\
            }
        }
    }
     \\ \\

\end{supertabular}
}

\end{document}
